\textbf{Case 2: Detecting a development bug in Conv2D semantics (CNN).}
During implementation of our Conv2D abstract transformer for the \textit{interval} domain, we encountered a subtle bug: the stride computation used \texttt{stride + 1} instead of \texttt{stride}, causing semantic drift between the verifier and the concrete model. Testing this implementation on a LeNet-style CNN under a seedable \textsf{BOX} specification, CBR returned \Inconclusive under the default 20-sample budget due to high input dimensionality. BBL, however, immediately detected the discrepancy: it flagged a large containment violation (gap $=0.45$) at the first Conv2D layer, ranking it as the top violation with downstream cascading effects. Guided by BBL's localization, we corrected the stride parameter. After the fix, BBL returned \Pass and the verifier's \Certified\ outcome correctly reflected the model's actual behavior. This real-world debugging scenario demonstrates BBL's practical value for catching implementation errors during verifier development.
